% Generated draft: requires mathematical and source review.
\documentclass[11pt]{article}
\usepackage[margin=25mm]{geometry}
\usepackage{fontspec}
\setmainfont{Latin Modern Roman}
\usepackage{amsmath,amssymb,mathtools,unicode-math}
\setmathfont{Latin Modern Math}
\usepackage{enumitem,needspace,xurl,hyperref}
\setlength{\parindent}{0pt}
\setlength{\parskip}{5pt}
\setlength{\emergencystretch}{3em}
\newcommand{\N}{\mathbb{N}}
\newcommand{\Z}{\mathbb{Z}}
\newcommand{\Q}{\mathbb{Q}}
\newcommand{\R}{\mathbb{R}}
\newcommand{\C}{\mathbb{C}}
\newcommand{\F}{\mathbb{F}}
\newcommand{\A}{\mathbb{A}}
\newcommand{\PP}{\mathbb{P}}
\newcommand{\E}{\mathbb{E}}
\newcommand{\eps}{\varepsilon}
\newcommand{\dd}{\,\mathrm{d}}
\begin{document}
\section*{UnsolvedMath ID: 20000326}
\textbf{Source problem code:} catalogue entry
\subsection{Definitions and mathematical statement}
The AIM question asks for a quasi-projectivity theorem intermediate between the result of Viehweg and the statement of Schumacher and Tsuji, for instance assuming the canonical class is effective. The record does not supply those comparison statements or the moduli functor, dimension, polarization, and boundedness hypotheses, so a unique general theorem cannot be reconstructed. Its subsequent paragraph notes Koll\'ar's point-blowup mechanism, parameterized by a variety modulo automorphisms, and speculates that every algebraic space expressible as a quotient of a quasi-projective variety by a proper group action might occur as a moduli space for a class of polarized manifolds. This is a suggested possibility, not an asserted realization theorem.

The report proposes a restricted fixed-parent result. Let \(Y\) be a fixed smooth projective complex surface with \(K_Y\) nef and linearly effective, let \(H\) be very ample, and fix an integer \(r\geq1\). For an ordered tuple \(P=(p_1,\ldots,p_r)\) of distinct points, let \(\pi_P:X_P\to Y\) be their blowup with exceptional curves \(E_1,\ldots,E_r\), and set
\[
L_P=\pi_P^*(H^{\otimes(r+1)})\otimes\mathcal O_{X_P}(-E_1-\cdots-E_r).
\]
The report asserts this gives polarized surfaces with
\[
K_{X_P}=\pi_P^*K_Y+\sum_iE_i
\]
linearly effective but non-nef, and that the irreducible \(K_{X_P}\)-negative curves are intrinsically exactly the exceptional curves. The non-nef assertion requires \(r\geq1\); the unblown case is separate.

It further asserts that every polarized isomorphism \((X_P,L_P)\cong(X_Q,L_Q)\) descends to an automorphism of \(Y\) preserving \(H^{\otimes(r+1)}\), carrying the unordered center sets to each other. The precise group is
\[
G=\operatorname{Aut}(Y,H^{\otimes(r+1)}),
\]
which the report says is finite under these hypotheses, rather than silently replacing it with \(\operatorname{Aut}(Y,H)\). Thus its proposed restricted coarse orbit space is
\[
\operatorname{Conf}_r(Y)/(G\times S_r),\qquad
\operatorname{Conf}_r(Y)=Y^r\setminus\bigcup_{i<j}\{p_i=p_j\},
\]
a quasi-projective finite quotient. The permutation group removes the ordering of centers. This is a coarse space for the indicated fixed-parent blowup family, not a claim that it is an entire component of an unrestricted polarized moduli space or a fine moduli scheme.

The report concludes that this simplest fixed-parent surface construction cannot produce an effective-canonical non-quasi-projective example. It does not settle the original general quasi-projectivity question, other parents or dimensions, or the speculative realization of arbitrary proper-action quotients.
\subsection{Short English statement}
Seek quasi-projectivity under a canonical-effectivity hypothesis. The report studies only blowups of a fixed minimal complex surface and proposes that their polarized coarse orbit space is a finite quasi-projective quotient, limiting that particular counterexample strategy.
\subsection{Sources}
\begin{itemize}
\item AIM workshop ``Compact moduli spaces and birational geometry,'' Problem 1.7: \url{https://aimath.org/WWN/birational/birational.pdf}. No PDF was accessed here.
\item Local canonical registry \texttt{aim-algebraic-geometry-notes.json}, \texttt{notes[325]}, tag \texttt{problem}; permanent UnsolvedMath ID \(20000326\), source code \texttt{catalogue entry}. Original names Viehweg, Schumacher, Tsuji, and Koll\'ar, and includes the speculative proper-action quotient paragraph. An interposed running header and page number are layout artifacts. Original remarks and literature field are empty.
\item Supplied research attempt 1 proposing the intrinsic \(K\)-negative recovery and finite-quotient result for simultaneous point blowups. No authoritative full bibliographic references or complete statements of the comparison quasi-projectivity results are supplied. The report is machine generated, not peer reviewed, and requires independent verification before citation.
\end{itemize}
\end{document}
